% ============================================================
% PAGE 5 — Proposed Method: Bidirectional Feature Pyramid Neck
% ============================================================

\subsection{Module 2: Bidirectional Attention-Augmented Neck}

The neck is the architectural module responsible for aggregating
the four backbone feature maps into a set of semantically
consistent, spatially precise feature representations suitable
for the detection head.  Standard FPN~\cite{lin2017feature}
accomplishes this via a single top-down pathway: beginning from
the deepest, most semantically rich feature map $P_5$, it
progressively upsamples and merges features into the shallower
levels $P_4$, $P_3$, $P_2$ via lateral connections.  Each
lateral connection adds the upsampled higher-level features to
the same-resolution lower-level features after a $1 \times 1$
convolution projects both to a common channel dimension.

Although effective, standard FPN has two well-documented
limitations.  First, it is \textit{unidirectional}: information
flows top-down only, so shallow feature maps at $P_2$ and $P_3$
benefit from high-level semantics but the updated shallow maps
never feed back to refine the deep maps.  PANet~\cite{liu2018path}
introduced a complementary bottom-up pathway that addresses
this limitation by propagating spatially precise low-level
information back up to the deeper levels.  Second, standard
lateral connections treat all feature channels as equally
important, whereas certain channels encode semantically
discriminative patterns that are more predictive of object
presence and should receive higher weight during fusion.

FusionDet's neck addresses both limitations through a
\textit{bidirectional fusion} topology augmented with
channel-and-spatial attention recalibration after each fusion
node.

\subsubsection{Bidirectional Fusion Topology}

The neck consists of two sequential passes over the four
pyramid levels:

\textbf{Top-down pass.}  Beginning from $P_5$, a $2\times$
bilinear upsampling operation halves the stride and doubles the
spatial dimensions, and the result is concatenated with $P_4$
along the channel axis.  A Fusion Block (two stacked
$3\times3$ depthwise-separable convolutions with batch
normalization and SiLU activation) processes the concatenated
tensor and produces an intermediate feature map $\tilde{P}_4$.
The process repeats for levels 3 and 2, producing $\tilde{P}_3$
and $\tilde{P}_2$.

\textbf{Bottom-up pass.}  Beginning from $\tilde{P}_2$, a
stride-2 depthwise-separable convolution downsamples the
feature map and the result is concatenated with $\tilde{P}_3$.
A second Fusion Block produces $\hat{P}_3$.  This continues
upward through $\hat{P}_4$ and $\hat{P}_5$, propagating the
spatially precise information that was enriched during the
top-down pass back to the deeper, more semantically abstracted
levels.

The complete sequence produces four recalibrated feature maps
$\hat{P}_2, \hat{P}_3, \hat{P}_4, \hat{P}_5$ that will serve
as inputs to the detection head.  The bidirectional topology
shortens the maximum information path between any two pyramid
levels from $O(L)$ (standard FPN) to $O(1)$ by ensuring that
every level directly receives both top-down and bottom-up
signals, where $L$ is the number of pyramid levels.

\subsubsection{Fusion Block Design}

Each Fusion Block consists of two successive $3 \times 3$
depthwise-separable convolutions:
\begin{equation}
    \mathbf{F}_{\text{out}} =
    \text{DWSConv}_{3\times3}\!\left(
        \text{DWSConv}_{3\times3}(\mathbf{F}_{\text{in}})
    \right)
    \label{eq:fblock}
\end{equation}
where $\text{DWSConv}_{3\times3}$ denotes a depthwise
convolution followed by a pointwise convolution, batch
normalization, and SiLU activation.  The two-layer structure
provides a receptive field of $5 \times 5$ pixels per fusion
node at marginal cost; replacing it with a single layer
degrades mAP by $0.6$ points in our ablation study (see
Table~\ref{tab:neck_ablation}).

\subsubsection{Channel-and-Spatial Attention Recalibration}

Following each Fusion Block, the feature map is passed through
a Convolutional Block Attention Module (CBAM)~\cite{woo2018cbam}
to produce a recalibrated feature map.

\textbf{Channel attention} captures inter-channel dependencies
by modeling which feature channels are most informative for
detecting objects at a given pyramid level.  Both global
average pooling and global max pooling summarize each channel
to a scalar value; both summaries are fed through a shared
two-layer MLP and summed:
\begin{equation}
    \mathbf{M}_c(\mathbf{F}) =
    \sigma\!\left(
        \text{MLP}\!\left(\text{AvgPool}(\mathbf{F})\right)
        + \text{MLP}\!\left(\text{MaxPool}(\mathbf{F})\right)
    \right)
    \label{eq:channel_att}
\end{equation}
where $\sigma$ denotes the sigmoid function.  The resulting
channel attention map $\mathbf{M}_c \in \mathbb{R}^{C \times
1 \times 1}$ is broadcast-multiplied with $\mathbf{F}$ to
produce a channel-recalibrated feature map
$\mathbf{F}' = \mathbf{M}_c(\mathbf{F}) \otimes \mathbf{F}$.

\textbf{Spatial attention} identifies which spatial locations
in the feature map are most informative, complementing the
channel-level recalibration.  Channel-wise average and max
pooling collapse the channel dimension to two scalar maps;
these are concatenated and processed by a $7 \times 7$
convolution:
\begin{equation}
    \mathbf{M}_s(\mathbf{F}') =
    \sigma\!\left(
        f^{7\times7}\!\left(
        [\text{AvgPool}_c(\mathbf{F}');\,
         \text{MaxPool}_c(\mathbf{F}')]
        \right)
    \right)
    \label{eq:spatial_att}
\end{equation}
where $[\cdot;\cdot]$ denotes concatenation along the channel
axis.  The spatial attention map $\mathbf{M}_s \in \mathbb{R}^{
1 \times H' \times W'}$ is broadcast-multiplied with
$\mathbf{F}'$ to produce the final recalibrated output:
\begin{equation}
    \hat{\mathbf{N}} =
    \mathbf{M}_s\!\left(\mathbf{F}'\right) \otimes \mathbf{F}'
    \label{eq:recalib}
\end{equation}

By explicitly directing both channel selection and spatial
focus, the CBAM recalibration encourages the neck to suppress
non-informative background activations and amplify object-relevant
feature responses before they reach the detection head.  This
is particularly beneficial for small objects, which occupy only
a small fraction of the feature map and are easily overwhelmed
by surrounding context; the spatial attention map learns to
concentrate gradient energy on those sparse locations.

\subsubsection{Neck Ablation}

Table~\ref{tab:neck_ablation} quantifies the individual
contribution of each neck component.

\begin{table}[!t]
\centering
\caption{Neck component ablation on MS-COCO val2017.
The baseline uses a standard FPN neck.}
\label{tab:neck_ablation}
\renewcommand{\arraystretch}{1.15}
\setlength{\tabcolsep}{3pt}
\begin{tabular}{l c c c}
\hline\hline
\textbf{Configuration}   & \textbf{mAP} &
\textbf{AP\textsubscript{S}} & \textbf{FPS} \\
\hline
FPN (baseline)           & 44.1 & 26.3 & 87 \\
+ Bottom-up path (PANet) & 45.0 & 27.4 & 82 \\
+ Two-layer Fusion Block & 45.6 & 28.1 & 80 \\
+ Channel attention      & 46.1 & 28.8 & 78 \\
\textbf{+ Spatial attention (full)} &
\textbf{46.7} & \textbf{29.5} & \textbf{76} \\
\hline\hline
\end{tabular}
\end{table}

Each component yields a monotonic mAP improvement, with the
largest single gain attributable to the bottom-up fusion path
($+0.9$ mAP) and the smallest to spatial attention alone
($+0.6$ mAP).  The FPS cost of the full neck versus the FPN
baseline is $11$ frames per second, a penalty judged acceptable
given the $2.6$-point aggregate mAP improvement.  Notably, the
small-object metric $\text{AP}_S$ improves by $3.2$ points,
confirming that the attention recalibration is particularly
effective for objects that occupy few feature-map pixels.
